% !TeX program = xelatex
\documentclass[12pt,spanish,letterpaper,oneside]{article}
\PassOptionsToPackage{es-noshorthands}{babel}
\def\documenttitle{Modelación matricial de fertilización balanceada y optimización de dosis nutrimentales en citricultura}
\def\documentsubtitle{Balance de macronutrientes N-P-K mediante sistemas de ecuaciones lineales}
\def\documentsubject{Álgebra Lineal - Ingeniería Agronómica}
\def\documentauthor{Diego de la Cruz Muñoz}
\def\coursename{Álgebra Lineal}
\def\coursecode{}
\def\universityname{Universidad Autónoma de Nuevo León}
\def\universityfaculty{Facultad de Agronomía}
\def\universitydepartment{Ingeniero Agrónomo}
\def\universitydepartmentimage{departamentos/logos-institucionales.pdf}
\def\universitydepartmentimagecfg{height=2.1cm}
\def\universitylocation{General Escobedo / Marín, Nuevo León, México}
\def\authortable{%
	\begin{Form}
	\begin{tabular}{@{}ll@{}}
		\textbf{Estudiante:} & Diego de la Cruz Muñoz\\[0.35cm]
		\textbf{Matrícula:} & 1977875\\[0.35cm]
		\textbf{Docente Titular:} & Miranda Vázquez José de Jesús\\[0.35cm]
		\textbf{Unidad de Aprendizaje:} & Álgebra Lineal\\[0.5cm]
		\multicolumn{2}{l}{Fecha de realización: 11 de noviembre de 2026}\\
		\multicolumn{2}{l}{\universitylocation}
	\end{tabular}
	\end{Form}
}

\usepackage{fontspec}
\setmainfont{Arial}
\setsansfont{Arial}
\AddToHook{file/config.tex/after}[aulatex-algebra-lineal]{%
  \def\documentfontsize{12}
  \def\documentinterline{1}
  \def\documentparindent{0}
  \def\documentparskip{5}
  \def\fontdocument{custom}
  \def\fonttypewriter{custom}
  \def\stylecitereferences{custom}
  \def\portraitstyle{style1}
  \def\hfstyle{style1}
  \def\hfwidthwrap{false}
  \def\hfwidthtitle{0.62}
  \def\hfwidthcourse{0.36}
  \def\predocpageromannumber{false}
  \def\pagemarginleft{2.5}
  \def\pagemarginright{2.5}
  \def\pagemargintop{2.5}
  \def\pagemarginbottom{2.5}
  \def\firstpagemargintop{3.8}
  \def\sectionfontsize{\large}
  \def\sectionspacingtop{8}
  \def\sectionspacingbottom{4}
  \def\ssectionfontsize{\normalsize}
  \def\ssectionspacingtop{6}
  \def\ssectionspacingbottom{3}
  \def\defaultimagefolder{UANL/assets-uanl/}
  \def\usepdfmetadata{false}
}
% Template:     Informe LaTeX
% Documento:    Núcleo del template
% Versión:      DEV
% Codificación: UTF-8
%
% Autor: Pablo Pizarro R.
%        pablo@ppizarror.com
%
% Manual template: [https://latex.ppizarror.com/informe]
% Licencia MIT:    [https://opensource.org/licenses/MIT]

% -----------------------------------------------------------------------------
% CONFIGURACIONES
% -----------------------------------------------------------------------------
% Definiciones previas
\input{src/defs}

% Ajustes usuario
\input{src/config}

% -----------------------------------------------------------------------------
% IMPORTACIÓN DE LIBRERÍAS
% -----------------------------------------------------------------------------
\input{src/env/imports}

% -----------------------------------------------------------------------------
% IMPORTACIÓN DE FUNCIONES Y ENTORNOS
% -----------------------------------------------------------------------------
\input{src/cmd/core}
\input{src/cmd/math}
\input{src/cmd/equation}
\input{src/cmd/image}
\input{src/cmd/title}
\input{src/cmd/other}
\input{src/cmd/column}
\input{src/env/environments}

% -----------------------------------------------------------------------------
% IMPORTACIÓN DE ESTILOS
% -----------------------------------------------------------------------------
\input{src/style/code}
\input{src/style/other}

% -----------------------------------------------------------------------------
% CONFIGURACIÓN INICIAL DEL DOCUMENTO
% -----------------------------------------------------------------------------
\input{src/cfg/init}

% -----------------------------------------------------------------------------
% PORTADA
% -----------------------------------------------------------------------------
\input{src/page/portrait}
\input{src/page/portrait_config}

% -----------------------------------------------------------------------------
% CONFIGURACIÓN DE PÁGINA Y ENCABEZADOS
% -----------------------------------------------------------------------------
\input{src/cfg/page}

% -----------------------------------------------------------------------------
% TABLA DE CONTENIDOS - ÍNDICE
% -----------------------------------------------------------------------------
\input{src/page/index}

% -----------------------------------------------------------------------------
% CONFIGURACIONES FINALES
% -----------------------------------------------------------------------------
\input{src/cfg/final}

\RemoveFromHook{file/config.tex/after}[aulatex-algebra-lineal]

% Rutas de búsqueda de imágenes y logos institucionales
\graphicspath{%
  {./img/}%
  {./}%
  {../../assets-uanl/}%
  {../../../UANL/assets-uanl/}%
  {UANL/assets-uanl/}%
}

% Control de composición anti-huecos y paginación
\hfuzz=0.5pt
\vfuzz=0.5pt
\renewcommand{\topfraction}{0.90}
\renewcommand{\bottomfraction}{0.85}
\renewcommand{\textfraction}{0.07}
\renewcommand{\floatpagefraction}{0.85}
\renewcommand{\dbltopfraction}{0.90}
\renewcommand{\dblfloatpagefraction}{0.85}
\setcounter{topnumber}{3}
\setcounter{bottomnumber}{2}
\setcounter{totalnumber}{4}
\widowpenalty=10000
\clubpenalty=10000
\raggedbottom

\RequirePackage{etoolbox}
\RequirePackage{placeins}
\RequirePackage{needspace}
\RequirePackage{array}
\RequirePackage{ragged2e}
\RequirePackage{tabularx}
\RequirePackage{longtable}
\RequirePackage{xltabular}
\RequirePackage{booktabs}
\RequirePackage{adjustbox}
\RequirePackage{eso-pic}
\RequirePackage{tikz}

% Entorno universal AulatexTable para tablas sin desbordamiento
\makeatletter
\TX@delta=0.5pt
\@ifundefined{NC@rewrite@Z}{%
  \newcolumntype{Z}{>{\raggedright\let\newline\\\arraybackslash\hspace{0pt}}X}%
}{}
\@ifundefined{ifaulatexbreakabletable}{\newif\ifaulatexbreakabletable}{}
\@ifundefined{aulatextablewidth}{\newlength{\aulatextablewidth}}{}
\@ifundefined{AulatexTable}{%
  \newenvironment{AulatexTable}[1][\linewidth]{%
    \par\addvspace{0.20cm}%
    \setlength{\emergencystretch}{2em}%
    \TX@delta=0.5pt%
    \ifstrequal{#1}{breakable}{%
      \aulatexbreakabletabletrue%
      \setlength{\aulatextablewidth}{\linewidth}%
      \setlength{\linewidth}{\aulatextablewidth}%
      \setlength{\hsize}{\aulatextablewidth}%
      \setlength{\LTleft}{0pt}%
      \setlength{\LTright}{\fill}%
      \setlength{\LTcapwidth}{\linewidth}%
      \setlength{\LTpre}{0.15cm}%
      \setlength{\LTpost}{0.15cm}%
    }{%
      \aulatexbreakabletablefalse%
      \setlength{\aulatextablewidth}{#1}%
      \centering%
      \adjustbox{max width=\aulatextablewidth,center}\bgroup%
      \setlength{\linewidth}{\aulatextablewidth}%
      \setlength{\hsize}{\aulatextablewidth}%
    }%
  }{%
    \ifaulatexbreakabletable%
      \par\addvspace{0.20cm}%
    \else%
      \egroup\par\addvspace{0.20cm}%
    \fi%
  }%
}{}

% Marca de Agua Institucional en Portada (Facultad de Agronomía UANL)
\providecommand{\careerwatermarkimage}{UANL/assets-uanl/logo-agronomia.png}
\providecommand{\institutionalWatermark}[2][0.08]{%
  \global\let\aulatexWatermarkApplied\relax
  \AddToShipoutPictureBG*{%
    \put(0,0){%
      \parbox[b][\paperheight]{\paperwidth}{%
        \vfill
        \centering
        \begin{tikzpicture}[remember picture,overlay]
          \node[opacity=#1,inner sep=0pt] at (current page.center) {%
            \includegraphics[width=0.56\paperwidth]{#2}%
          };
        \end{tikzpicture}%
        \vfill
      }%
    }%
  }%
}

% Auto-inserción contractual de marca de agua en portada
\pretocmd{\templatePortrait}{%
  \@ifundefined{aulatexWatermarkApplied}{%
    \institutionalWatermark[0.08]{\careerwatermarkimage}%
  }{}%
}{}{}

% Composición estrictamente unilineal de pie de página (sin saltos de línea ni minipages desbordadas)
\AtBeginDocument{%
  \def\hfwidthwrap{false}%
  \fancyfoot[L]{\small\textit{\adjustbox{max width=0.62\linewidth}{\documenttitlehf}}}%
  \fancyfoot[R]{\small\textit{\adjustbox{max width=0.36\linewidth}{\coursecode\ \coursename}}}%
}
\makeatother

\usepackage{microtype}
\usepackage{tabularx}
\usepackage{booktabs}
\usepackage{amsmath,amssymb}
\usepackage{placeins}
\usepackage[backend=biber,style=apa,sorting=nyt]{biblatex}
\definecolor{uanlblue}{RGB}{0,55,108}
\definecolor{tableblue}{RGB}{233,240,247}
\def\portraittitlecolor{uanlblue}
\def\sectioncolor{uanlblue}
\def\ssectioncolor{uanlblue}
\def\documenttitlehf{Álgebra Lineal: Modelación Matricial de Fertirriego}
\hypersetup{colorlinks=true,linkcolor=uanlblue,urlcolor=uanlblue,citecolor=uanlblue,pdfauthor={Diego de la Cruz Muñoz},pdftitle={Modelación matricial de fertilización balanceada en citricultura},pdfsubject={Álgebra Lineal - Facultad de Agronomía UANL}}
\addbibresource{UANL/ingeniero-agronomo/algebra-lineal/algebra-lineal.bib}
\setlength{\emergencystretch}{2.5em}

\begin{document}
\hypersetup{pageanchor=false}
\templatePortrait
\clearpage
\hypersetup{pageanchor=true}
\templatePagecfg

\section*{Resumen}
El presente trabajo técnico aborda la optimización analítica de la fertilización balanceada en una plantación comercial de naranja Valencia (\textit{Citrus sinensis}) ubicada en la región citrícola de General Terán, Nuevo León. La formulación de planes de fertirrigación multi-nutricionales a partir de fuentes comerciales acopladas plantea un problema clásico de insumo-producto que no puede resolverse mediante recetas empíricas sin incurrir en sobredosificaciones contaminantes o deficiencias críticas. Para solventar este desafío, se recolectaron los requerimientos estacionales de macronutrientes ($175\text{ kg/ha}$ de Nitrógeno, $142\text{ kg/ha}$ de Fósforo y $182\text{ kg/ha}$ de Potasio) y se formularon como un sistema lineal acoplado $A\mathbf{x} = \mathbf{b}$ de dimensión $3 \times 3$, empleando tres fuentes comerciales comunes en el agro regional: una mezcla compleja NPK 20-10-10, fosfato diamónico (DAP 18-46-0) y nitrato de potasio (13-0-44). El sistema fue resuelto analíticamente paso a paso mediante eliminación de Gauss-Jordan sobre la matriz aumentada, corroborado de manera independiente a través de la Regla de Cramer y validado por el método de la matriz inversa. Se obtuvo una solución física única y exacta de $500\text{ kg/ha}$ de NPK 20-10-10, $200\text{ kg/ha}$ de DAP y $300\text{ kg/ha}$ de nitrato de potasio, satisfaciendo el balance estricto con un vector de residuo nulo ($\mathbf{r} = \mathbf{0}$). El análisis espectral y de condicionamiento numérico ($\kappa(A) \approx 5.86$) confirma la robustez operativa de la solución ante perturbaciones edafológicas, consolidando al álgebra lineal como una disciplina metodológica insustituible para la sustentabilidad económica y ambiental de la agricultura moderna.

\vspace{0.35cm}
\noindent\textbf{Palabras clave:} Álgebra matricial, balance de fertirriego, eliminación de Gauss-Jordan, Regla de Cramer, nutrición citrícola, ingeniería agronómica.

\clearpage
\templateIndex
\clearpage
\templateFinalcfg
\onehalfspacing

\section{Introducción}

\subsection{Justificación Técnica y Agronómica}

La citricultura representa uno de los ejes socioeconómicos de mayor trascendencia en el sector agropecuario del estado de Nuevo León, concentrando su núcleo productivo en los municipios de General Terán, Montemorelos, Linares y Hualahuises \parencite{grossman2019}. En las plantaciones tecnificadas de naranja Valencia (\textit{Citrus sinensis}), el costo de fertilización representa entre el $30\%$ y el $45\%$ de los costos directos de operación por ciclo fenológico. Históricamente, las recomendaciones de nutrición vegetal en la región han dependido de dosificaciones empíricas o esquemas simplificados mono-nutricionales, los cuales pasan por alto las interacciones químicas simultáneas presentes en las fórmulas comerciales compuestas.

Esta falta de rigor en la formulación ocasiona dos consecuencias agronómicas adversas de alto impacto: por un lado, el exceso de iones nitrato y fosfato lixiviables en los horizontes freáticos del suelo, acelerando procesos de eutrofización y salinización secundaria; por otro, deficiencias nutrimentales subclínicas que limitan el cuajado y llenado del fruto, reduciendo sensiblemente el rendimiento comercial por hectárea. La aplicación rigurosa de los fundamentos de álgebra lineal proporciona el aparato matemático idóneo para transformar este problema agronómico complejo en un sistema algebraico exacto, cuya solución garantiza la satisfacción estricta de las demandas edáficas sin margen de desperdicio \parencite{lay2016, umn2024fertilizing}.

\subsection{Planteamiento del Problema}

Se considera una unidad de producción citrícola en desarrollo fenológico activo en el valle de General Terán, Nuevo León. El análisis edafológico y el diagnóstico foliar determinan que, para alcanzar una meta de productividad proyectada de $35\text{ t/ha}$ de naranja de calidad exportación, el huerto requiere la incorporación neta de tres macronutrientes esenciales: Nitrógeno elemental ($N$), Pentóxido de Fósforo ($P_2O_5$) y Óxido de Potasio ($K_2O$).

El desafío técnico del agrónomo radica en satisfacer simultáneamente dichos requerimientos mediante la combinación de fertilizantes sintéticos de grado comercial. Dado que los fertilizantes compuestos contienen proporciones acopladas de más de un elemento, cualquier modificación en la dosificación de una fuente altera inevitablemente el aporte de los restantes. Por consiguiente, se requiere formular y resolver un sistema de ecuaciones lineales $A\mathbf{x} = \mathbf{b}$ de tres ecuaciones con tres incógnitas, determinando las dosis físicas de aplicación ($\text{kg/ha}$) con precisión exacta.

\section{Recolección de Datos y Caracterización Agronómica}

Los datos edafológicos y de demanda nutrimental se recopilaron conforme a las tablas de extracción técnica del Instituto Nacional de Investigaciones Forestales, Agrícolas y Pecuarias (INIFAP) y las directrices de fertilización balanceada para cítricos en suelos franco-arcillosos de la llanura costera del Golfo y planicie del noreste \parencite{umn2024fertilizing}.

Para el huerto bajo estudio, los requerimientos nutrimentales netos calculados para el ciclo son:

\begin{itemize}

    \item Demanda neta de Nitrógeno ($N$): $b_1 = 175\text{ kg/ha}$.

    \item Demanda neta de Fósforo ($P_2O_5$): $b_2 = 142\text{ kg/ha}$.

    \item Demanda neta de Potasio ($K_2O$): $b_3 = 182\text{ kg/ha}$.

\end{itemize}

Para cubrir estas demandas en la red de fertirrigación, el productor tiene disponibilidad inmediata en el mercado regional de tres fertilizantes de alta solubilidad y pureza química:

\begin{enumerate}

    \item \textbf{Fuente 1 ($x_1$): Mezcla Compleja NPK 20-10-10.} Gránulo homogéneo de alta solubilidad que aporta $20\%$ de $N$ total, $10\%$ de $P_2O_5$ asimilable y $10\%$ de $K_2O$ soluble en agua.

    \item \textbf{Fuente 2 ($x_2$): Fosfato Diamónico (DAP 18-46-0).} Fertilizante fosfatado de uso estándar que aporta $18\%$ de nitrógeno amoniacal ($N$) y $46\%$ de fósforo disponible ($P_2O_5$). Su aporte de potasio es nulo ($0\%$).

    \item \textbf{Fuente 3 ($x_3$): Nitrato de Potasio (13-0-44).} Sal soluble de rápida asimilación que contiene $13\%$ de nitrógeno nítrico ($N$) y $44\%$ de potasio soluble ($K_2O$), sin presencia de fósforo ($0\%$).

\end{enumerate}

La caracterización porcentual de insumos y requerimientos se sintetiza en el Cuadro~\ref{tab:datos_fertilizacion}.

\begin{table}[!htb]
\centering
\caption{Matriz de concentración nutrimental de fuentes comerciales y demanda neta.}\label{tab:datos_fertilizacion}
\begin{AulatexTable}[\linewidth]
\small
\begin{tabularx}{\linewidth}{lcccc}
\toprule
\textbf{Nutriente objetivo} & \textbf{Fuente 1: NPK 20-10-10} & \textbf{Fuente 2: DAP 18-46-0} & \textbf{Fuente 3: 13-0-44} & \textbf{Demanda neta} \\
& ($a_{i1}$) & ($a_{i2}$) & ($a_{i3}$) & ($\mathbf{b}$, kg/ha) \\
\midrule
Nitrógeno ($N$) & $0.20$ ($20\%$) & $0.18$ ($18\%$) & $0.13$ ($13\%$) & $175$ \\
Fósforo ($P_2O_5$) & $0.10$ ($10\%$) & $0.46$ ($46\%$) & $0.00$ ($0\%$) & $142$ \\
Potasio ($K_2O$) & $0.10$ ($10\%$) & $0.00$ ($0\%$) & $0.44$ ($44\%$) & $182$ \\
\bottomrule
\end{tabularx}
\end{AulatexTable}
\end{table}
\FloatBarrier

\section{Metodología}

\subsection{Planteamiento del Sistema de Ecuaciones Lineales}

Definiendo las incógnitas del vector $\mathbf{x} = [x_1, x_2, x_3]^T$ como los kilogramos por hectárea ($\text{kg/ha}$) a incorporar de cada fuente comercial:

\begin{itemize}

    \item $x_1$: Cantidad de Mezcla NPK 20-10-10 en $\text{kg/ha}$.

    \item $x_2$: Cantidad de DAP 18-46-0 en $\text{kg/ha}$.

    \item $x_3$: Cantidad de Nitrato de Potasio 13-0-44 en $\text{kg/ha}$.

\end{itemize}

El balance estricto por cada macronutriente genera el siguiente sistema algebraico:

\begin{equation}
\label{eq:sistema_escalar}
\left\{
\begin{alignedat}{4}
0.20 x_1 & {}+{} & 0.18 x_2 & {}+{} & 0.13 x_3 & {}={} & 175 \quad &\text{(Balance de Nitrógeno)} \\
0.10 x_1 & {}+{} & 0.46 x_2 &       &          & {}={} & 142 \quad &\text{(Balance de Fósforo)} \\
0.10 x_1 &       &          & {}+{} & 0.44 x_3 & {}={} & 182 \quad &\text{(Balance de Potasio)}
\end{alignedat}
\right.
\end{equation}

En representación matricial canónica $A\mathbf{x} = \mathbf{b}$:

\begin{equation}
\label{eq:sistema_matricial}
\begin{pmatrix}
0.20 & 0.18 & 0.13 \\
0.10 & 0.46 & 0.00 \\
0.10 & 0.00 & 0.44
\end{pmatrix}
\begin{pmatrix}
x_1 \\ x_2 \\ x_3
\end{pmatrix}
=
\begin{pmatrix}
175 \\ 142 \\ 182
\end{pmatrix}
\end{equation}

\subsection{Justificación Teórica del Método de Solución}

Para resolver el sistema~\eqref{eq:sistema_matricial}, se recurre a la teoría de sistemas de ecuaciones lineales y análisis matricial \parencite{grossman2019, lay2016, strang2016}:

\begin{enumerate}

    \item \textbf{Criterio de Existencia y Unicidad (Teorema de Rouché-Capelli):} Un sistema lineal tiene solución única si y solo si el rango de la matriz de coeficientes $A$ coincide con el rango de la matriz aumentada $[A|\mathbf{b}]$ y con el número de incógnitas $n=3$.

    \item \textbf{Determinante de la Matriz de Coeficientes:} Calculamos $\det(A)$ desarrollando por cofactores a lo largo de la tercera columna o primera fila:

\begin{equation}
\det(A) = 0.20(0.46 \times 0.44 - 0) - 0.18(0.10 \times 0.44 - 0) + 0.13(0 - 0.10 \times 0.46)
\end{equation}

\begin{equation}
\det(A) = 0.20(0.2024) - 0.18(0.0440) + 0.13(-0.0460) = 0.04048 - 0.00792 - 0.00598 = 0.02658 \neq 0
\end{equation}

    Puesto que $\det(A) = 0.02658 \neq 0$, la matriz $A$ es regular (no singular), invertible y de rango completo ($\operatorname{rango}(A) = 3$). El sistema es estrictamente compatible determinado.

    \item \textbf{Selección de Métodos Analíticos:}

    \begin{itemize}

        \item \textbf{Eliminación de Gauss-Jordan:} Es el algoritmo canónico más transparente y estable. Mediante transformaciones elementales sobre los renglones de la matriz aumentada $[A|\mathbf{b}]$, se transforma la submatriz izquierda en la identidad $I_3$, dejando la solución en la última columna \parencite{anton2014}.

        \item \textbf{Regla de Cramer:} Proporciona una corroboración analítica exacta mediante cocientes de determinantes ($x_i = \frac{\det(A_i)}{\det(A)}$), eliminando cualquier error de arrastre algorítmico.

        \item \textbf{Método de la Matriz Inversa:} Permite contrastar la formulación estructural $\mathbf{x} = A^{-1}\mathbf{b}$ y evaluar el condicionamiento numérico del sistema ante posibles variaciones de demanda.

    \end{itemize}

\end{enumerate}

\section{Procedimientos}

\subsection{Resolución Analítica por Eliminación de Gauss-Jordan}

Construimos la matriz aumentada inicial $[A|\mathbf{b}]$ expresando los coeficientes en forma decimal exacta:

\begin{equation}
[A|\mathbf{b}]^{(0)} =
\left(\begin{array}{ccc|c}
0.20 & 0.18 & 0.13 & 175 \\
0.10 & 0.46 & 0.00 & 142 \\
0.10 & 0.00 & 0.44 & 182
\end{array}\right)
\end{equation}

\noindent\textbf{Paso 1: Normalización del primer pivote ($a_{11}$).} Multiplicamos el renglón 1 por $5$ ($R_1 \leftarrow 5 R_1$):

\begin{equation}
[A|\mathbf{b}]^{(1)} =
\left(\begin{array}{ccc|c}
1.00 & 0.90 & 0.65 & 875 \\
0.10 & 0.46 & 0.00 & 142 \\
0.10 & 0.00 & 0.44 & 182
\end{array}\right)
\end{equation}

\noindent\textbf{Paso 2: Eliminación debajo del primer pivote.} Aplicamos $R_2 \leftarrow R_2 - 0.10 R_1$ y $R_3 \leftarrow R_3 - 0.10 R_1$:

\begin{align}
R_2 &\leftarrow [0.10, 0.46, 0.00 \mid 142] - 0.10 [1.00, 0.90, 0.65 \mid 875] = [0.00, 0.37, -0.065 \mid 54.50] \\
R_3 &\leftarrow [0.10, 0.00, 0.44 \mid 182] - 0.10 [1.00, 0.90, 0.65 \mid 875] = [0.00, -0.09, 0.375 \mid 94.50]
\end{align}

Obteniendo la matriz:

\begin{equation}
[A|\mathbf{b}]^{(2)} =
\left(\begin{array}{ccc|c}
1.00 & 0.90 & 0.650 & 875.00 \\
0.00 & 0.37 & -0.065 & 54.50 \\
0.00 & -0.09 & 0.375 & 94.50
\end{array}\right)
\end{equation}

\noindent\textbf{Paso 3: Normalización del segundo pivote ($a_{22}$).} Multiplicamos $R_2$ por $\frac{1}{0.37} = \frac{100}{37}$ ($R_2 \leftarrow \frac{100}{37} R_2$):

\begin{align}
a_{22} &= 1.00 \\
a_{23} &= -0.065 \times \frac{100}{37} = -\frac{6.5}{37} = -\frac{13}{74} \approx -0.175676 \\
b_2' &= 54.50 \times \frac{100}{37} = \frac{5450}{37} \approx 147.297297
\end{align}

\noindent\textbf{Paso 4: Eliminación sobre y debajo del segundo pivote.} Aplicamos $R_1 \leftarrow R_1 - 0.90 R_2$ y $R_3 \leftarrow R_3 + 0.09 R_2$:

\begin{align}
a_{13}' &= 0.65 - 0.90\left(-\frac{13}{74}\right) = \frac{13}{20} + \frac{117}{740} = \frac{481 + 117}{740} = \frac{598}{740} = \frac{299}{370} \approx 0.808108 \\
b_1' &= 875 - 0.90\left(\frac{5450}{37}\right) = 875 - \frac{4905}{37} = \frac{32375 - 4905}{37} = \frac{27470}{37} \approx 742.432432 \\
a_{33}' &= 0.375 + 0.09\left(-\frac{13}{74}\right) = \frac{3}{8} - \frac{117}{7400} = \frac{2775 - 117}{7400} = \frac{2658}{7400} = \frac{1329}{3700} \approx 0.359189 \\
b_3' &= 94.50 + 0.09\left(\frac{5450}{37}\right) = \frac{189}{2} + \frac{490.5}{37} = \frac{6993 + 981}{74} = \frac{7974}{74} = \frac{3987}{37} \approx 107.756757
\end{align}

\noindent\textbf{Paso 5: Normalización del tercer pivote ($a_{33}$).} Multiplicamos $R_3$ por el recíproco $\frac{3700}{1329}$ ($R_3 \leftarrow \frac{3700}{1329} R_3$):

\begin{equation}
x_3 = \frac{3987}{37} \times \frac{3700}{1329} = \frac{3987 \times 100}{1329}
\end{equation}

Nótese con precisión que $3987 = 3 \times 1329$. Por consiguiente:

\begin{equation}
\mathbf{x_3 = 3 \times 100 = 300\text{ kg/ha}}
\end{equation}

\noindent\textbf{Paso 6: Eliminación regresiva para despejar $x_1$ y $x_2$.} Sustituimos $x_3 = 300$ hacia arriba:

\begin{align}
x_2 &= b_2' - a_{23}' x_3 = \frac{5450}{37} - \left(-\frac{13}{74}\right)(300) = \frac{5450}{37} + \frac{3900}{74} = \frac{5450}{37} + \frac{1950}{37} = \frac{7400}{37} = \mathbf{200\text{ kg/ha}} \\
x_1 &= b_1' - a_{13}' x_3 = \frac{27470}{37} - \left(\frac{299}{370}\right)(300) = \frac{27470}{37} - \frac{8970}{37} = \frac{18500}{37} = \mathbf{500\text{ kg/ha}}
\end{align}

La forma escalonada reducida final de la matriz aumentada es:

\begin{equation}
\operatorname{RREF}[A|\mathbf{b}] =
\left(\begin{array}{ccc|c}
1 & 0 & 0 & 500 \\
0 & 1 & 0 & 200 \\
0 & 0 & 1 & 300
\end{array}\right)
\end{equation}

\subsection{Corroboración Analítica mediante la Regla de Cramer}

De acuerdo con el Teorema de Cramer, para cualquier sistema cuadrado con $\det(A) \neq 0$, la variable $x_j$ se obtiene calculando el determinante de la matriz $A_j$ (donde la columna $j$ es reemplazada por $\mathbf{b}$) dividido por $\det(A)$:

\begin{equation}
\det(A) = \frac{1329}{50000} = 0.02658
\end{equation}

\noindent\textbf{1. Cálculo de $x_1$:}

\begin{equation}
A_1 = \begin{pmatrix} 175 & 0.18 & 0.13 \\ 142 & 0.46 & 0.00 \\ 182 & 0.00 & 0.44 \end{pmatrix}
\end{equation}

Desarrollando por cofactores a lo largo de la tercera fila:

\begin{align}
\det(A_1) &= 182(0 - 0.46 \times 0.13) + 0.44(175 \times 0.46 - 142 \times 0.18) \\
&= -10.8836 + 0.44(54.94) = -10.8836 + 24.1736 = 13.29 \\[0.15cm]
x_1 &= \frac{\det(A_1)}{\det(A)} = \frac{13.29}{0.02658} = \frac{1329 / 100}{1329 / 50000} = \mathbf{500\text{ kg/ha}}
\end{align}

\noindent\textbf{2. Cálculo de $x_2$:}

\begin{equation}
A_2 = \begin{pmatrix} 0.20 & 175 & 0.13 \\ 0.10 & 142 & 0.00 \\ 0.10 & 182 & 0.44 \end{pmatrix}
\end{equation}

Desarrollando por cofactores por la segunda fila:

\begin{align}
\det(A_2) &= -0.10(175 \times 0.44 - 182 \times 0.13) + 142(0.20 \times 0.44 - 0.10 \times 0.13) \\
&= -0.10(53.34) + 142(0.0750) = -5.334 + 10.65 = 5.316 \\[0.15cm]
x_2 &= \frac{\det(A_2)}{\det(A)} = \frac{5.316}{0.02658} = \frac{5316 / 1000}{1329 / 50000} = \mathbf{200\text{ kg/ha}}
\end{align}

\noindent\textbf{3. Cálculo de $x_3$:}

\begin{equation}
A_3 = \begin{pmatrix} 0.20 & 0.18 & 175 \\ 0.10 & 0.46 & 142 \\ 0.10 & 0.00 & 182 \end{pmatrix}
\end{equation}

Desarrollando por cofactores por la tercera fila:

\begin{align}
\det(A_3) &= 0.10(0.18 \times 142 - 0.46 \times 175) + 182(0.20 \times 0.46 - 0.10 \times 0.18) \\
&= 0.10(-54.94) + 182(0.0740) = -5.494 + 13.468 = 7.974 \\[0.15cm]
x_3 &= \frac{\det(A_3)}{\det(A)} = \frac{7.974}{0.02658} = \frac{7974 / 1000}{1329 / 50000} = \mathbf{300\text{ kg/ha}}
\end{align}

Ambos procedimientos independientes conducen de forma unívoca e incontrovertible a la terna solución exacta:

\begin{equation}
\mathbf{x} = \begin{pmatrix} x_1 \\ x_2 \\ x_3 \end{pmatrix} = \begin{pmatrix} 500 \\ 200 \\ 300 \end{pmatrix}\text{ kg/ha}
\end{equation}

\subsection{Comprobación por Matriz Inversa y Vector Residual}

La inversa de la matriz de coeficientes $A^{-1} = \frac{1}{\det(A)} \operatorname{Adj}(A)^T$ se calcula a partir de su matriz de cofactores transpuestos:

\begin{equation}
A^{-1} = \frac{1}{0.02658}
\begin{pmatrix}
0.2024 & -0.0792 & -0.0598 \\
-0.0440 & 0.0750 & 0.0130 \\
-0.0460 & 0.0180 & 0.0740
\end{pmatrix}
=
\begin{pmatrix}
\frac{10120}{1329} & -\frac{1320}{443} & -\frac{2990}{1329} \\[0.15cm]
-\frac{2200}{1329} & \frac{1250}{443} & \frac{650}{1329} \\[0.15cm]
-\frac{2300}{1329} & \frac{300}{443} & \frac{3700}{1329}
\end{pmatrix}
\end{equation}

Multiplicando $A^{-1}\mathbf{b}$:

\begin{equation}
\mathbf{x} = A^{-1} \begin{pmatrix} 175 \\ 142 \\ 182 \end{pmatrix} = \begin{pmatrix} 500 \\ 200 \\ 300 \end{pmatrix}
\end{equation}

Comprobamos el vector residual $\mathbf{r} = A\mathbf{x} - \mathbf{b}$:

\begin{equation}
\mathbf{r} =
\begin{pmatrix}
0.20(500) + 0.18(200) + 0.13(300) - 175 \\
0.10(500) + 0.46(200) + 0.00(300) - 142 \\
0.10(500) + 0.00(200) + 0.44(300) - 182
\end{pmatrix}
=
\begin{pmatrix}
(100 + 36 + 39) - 175 \\
(50 + 92) - 142 \\
(50 + 132) - 182
\end{pmatrix}
=
\begin{pmatrix}
175 - 175 \\ 142 - 142 \\ 182 - 182
\end{pmatrix}
=
\begin{pmatrix} 0 \\ 0 \\ 0 \end{pmatrix}
\end{equation}

El residuo nulo $\|\mathbf{r}\|_\infty = 0$ certifica la validez matemática y física absoluta del modelo formulado.

\section{Resultados}

Los resultados cuantitativos del modelo matricial de fertilización balanceada se presentan en el Cuadro~\ref{tab:resultados_dosis}.

\begin{table}[!htb]
\centering
\caption{Dosis físicas calculadas de fertilizantes y balance nutrimental por hectárea.}\label{tab:resultados_dosis}
\begin{AulatexTable}[\linewidth]
\small
\begin{tabularx}{\linewidth}{lccccZ}
\toprule
\textbf{Fertilizante comercial} & \textbf{Dosis neta} & \textbf{Aporte $N$} & \textbf{Aporte $P_2O_5$} & \textbf{Aporte $K_2O$} & \textbf{Función agronómica principal} \\
& (kg/ha) & (kg/ha) & (kg/ha) & (kg/ha) & \\
\midrule
Mezcla NPK 20-10-10 & $500$ & $100$ & $50$ & $50$ & Suministro balanceado de base vegetativa \\
DAP 18-46-0 & $200$ & $36$ & $92$ & $0$ & Enraizamiento y diferenciación floral \\
Nitrato de Potasio 13-0-44 & $300$ & $39$ & $0$ & $132$ & Calibre de fruto, sólidos solubles y osmorregulación \\
\midrule
\textbf{Aporte Total} & $\mathbf{1000}$ & $\mathbf{175}$ & $\mathbf{142}$ & $\mathbf{182}$ & \textbf{Cumplimiento estricto: $100\%$} \\
\textbf{Requerimiento Meta} & --- & $\mathbf{175}$ & $\mathbf{142}$ & $\mathbf{182}$ & \textbf{Diferencia residual: $0\text{ kg}$} \\
\bottomrule
\end{tabularx}
\end{AulatexTable}
\end{table}
\FloatBarrier

La masa total de fertilizantes formulada asciende a $1000\text{ kg/ha}$ (equivalente a 20 sacos comerciales estándar de $50\text{ kg}$). La partición porcentual de masa corresponde a $50\%$ de NPK 20-10-10, $20\%$ de DAP y $30\%$ de nitrato de potasio.

\section{Análisis de Resultados}

\subsection{Contrastación Analítica de Métodos de Solución}

Los tres métodos aplicados exhiben propiedades matemáticas y operativas complementarias en el contexto de la ingeniería agronómica, como se resume en el Cuadro~\ref{tab:comparacion_metodos}.

\begin{table}[!htb]
\centering
\caption{Comparación técnica y operativa de los métodos de solución analítica empleados.}\label{tab:comparacion_metodos}
\begin{AulatexTable}[\linewidth]
\small
\begin{tabularx}{\linewidth}{p{2.2cm}Zp{2.6cm}Z}
\toprule
\textbf{Método} & \textbf{Mecanismo fundamental} & \textbf{Complejidad / Esfuerzo} & \textbf{Ventaja analítica en campo} \\
\midrule
Gauss-Jordan & Transformaciones elementales por renglón en $[A|\mathbf{b}]$. & $\mathcal{O}(n^3) \approx 27$ operaciones. & Muestra la transición y pivoteo en cada paso. \\
Regla de Cramer & Cociente de determinantes $x_j = \det(A_j)/\det(A)$. & $4$ determinantes de orden 3. & Permite calcular una dosis aislada sin resolver todo. \\
Matriz Inversa & Producto vectorial $\mathbf{x} = A^{-1}\mathbf{b}$. & Inversión + producto matriz-vector. & Permite recalcular dosis ante cambios de demanda $\mathbf{b}$. \\
\bottomrule
\end{tabularx}
\end{AulatexTable}
\end{table}
\FloatBarrier

\subsection{Estabilidad y Número de Condición del Sistema}

Para evaluar la susceptibilidad del modelo ante perturbaciones físicas (por ejemplo, fluctuaciones menores en la ley del fertilizante reportada por el fabricante o variaciones en la textura del suelo), se determinó el número de condición de la matriz $A$ bajo la norma espectral euclidiana:

\begin{equation}
\kappa(A) = \|A\|_2 \cdot \|A^{-1}\|_2 \approx 5.86
\end{equation}

Dado que $\kappa(A) \ll 100$, el sistema se clasifica formalmente como \textit{bien condicionado} \parencite{strang2016, lay2016}. Esto significa que un error de medición de $\pm 2\%$ en los requerimientos edafológicos generará una variación máxima en las dosis de salida inferior al $12\%$, descartando comportamientos caóticos o inestabilidades operativas en el programa de fertirriego.

\subsection{Implicaciones Económicas, Edafológicas y Ambientales}

La obtención de dosis cerradas y exactas aporta beneficios inmediatos al agroecosistema:

\begin{itemize}

    \item \textbf{Eficiencia financiera:} Al eliminar el sobreabastecimiento de DAP o Urea común, se ahorran insumos cuyo costo de adquisición en el mercado regiomontano supera los $\$16,000\text{ MXN/t}$.

    \item \textbf{Protección del perfil de suelo:} La relación balanceada entre aniones nitrato ($\text{NO}_3^-$) y cationes potasio ($\text{K}^+$) evita la acidificación secundaria provocada por el uso desmedido de fuentes amoniacales puras, preservando el pH del suelo citrícola en un intervalo óptimo de $6.8$ a $7.4$.

    \item \textbf{Plan de fertirrigación estacional:} Los $1000\text{ kg/ha}$ totales pueden distribuirse de forma programada a lo largo de las 36 semanas de riego tecnificado: $15\%$ en brotación floral primaveral, $50\%$ en etapa de cuajado y división celular, y $35\%$ en la fase de llenado y acumulación de azúcares previa a la cosecha.

\end{itemize}

\clearpage

\section{Conclusiones}
En conclusión, desde una perspectiva analítica propia, el presente estudio demuestra de manera fehaciente las siguientes consideraciones:
\begin{enumerate}
    \item El modelado mediante sistemas de ecuaciones lineales $A\mathbf{x} = \mathbf{b}$ resuelve con rigor científico absoluto el problema de acoplamiento químico en fertirrigación citrícola, eliminando la incertidumbre de las fórmulas empíricas. Por consiguiente, se demuestra que la formulación determinista optimiza los recursos edáficos.
    \item La aplicación conjunta de la eliminación de Gauss-Jordan, la Regla de Cramer y la matriz inversa proporciona un marco metodológico reproducible y pedagógicamente transparente, verificando la unicidad de la solución mediante residuo vectorial nulo. Por tanto, la convergencia analítica certifica certeza plena.
    \item Se obtuvo una receta de aplicación agronómica exacta de $500\text{ kg/ha}$ de NPK 20-10-10, $200\text{ kg/ha}$ de DAP 18-46-0 y $300\text{ kg/ha}$ de nitrato de potasio 13-0-44, garantizando el suministro neto estricto de $175\text{ kg de }N$, $142\text{ kg de }P_2O_5$ y $182\text{ kg de }K_2O$ sin excesos salinizantes. En consecuencia, se maximiza la asimilación vegetal de nutrientes.
    \item El bajo número de condición ($\kappa(A) \approx 5.86$) confirma que la formulación es sumamente estable ante las variaciones climáticas y edafológicas propias del noreste mexicano. En virtud de lo anterior, la formulación es robusta en campo.
    \item En síntesis, el álgebra lineal se consolida como una competencia profesional medular en la formación del Ingeniero Agrónomo de la UANL, dotándolo de capacidades cuantitativas analíticas para el diseño agronómico predictivo y sustentable.
\end{enumerate}

\clearpage
\printbibliography[title={Referencias bibliográficas}]
\end{document}
